\documentclass[12pt,oneside]{book}

% ============================================================
% METADATOS
% ============================================================
\newcommand{\universidad}{Universidad de Buenos Aires (UBA)}
\newcommand{\carrera}{Carrera de Abogacía}
\newcommand{\materia}{Lectocomprensión}
\newcommand{\titulo}{Inglés jurídico: contratos y lectura de fallos}
\newcommand{\autor}{Isaac Edgar Camacho Ocampo}
\newcommand{\anio}{2026}

% ============================================================
% PAQUETES
% ============================================================
\usepackage[utf8]{inputenc}
\usepackage[T1]{fontenc}
\usepackage[spanish,es-nodecimaldot]{babel}

\usepackage[
    top=2.5cm,
    bottom=2.5cm,
    left=3cm,
    right=2.5cm,
    headheight=15pt
]{geometry}

\usepackage{amsmath}
\usepackage{amssymb}
\usepackage{mathtools}

\usepackage{graphicx}
\usepackage{float}
\usepackage{caption}
\usepackage{subcaption}

\usepackage{booktabs}
\usepackage{array}
\usepackage{longtable}
\usepackage{tabularx}

\usepackage{enumitem}

\usepackage{xcolor}
\usepackage[most]{tcolorbox}

\usepackage{fancyhdr}

\usepackage{ragged2e}

% ============================================================
% RUTAS DE IMÁGENES
% ============================================================
\graphicspath{
    {./img/}
    {./graficos/}
    {./figuras/}
}

% ============================================================
% CONFIGURACIÓN GENERAL
% ============================================================
\setcounter{tocdepth}{2}

\setlength{\parindent}{0pt}
\setlength{\parskip}{0.5em}

\setlist[itemize]{
    topsep=0.3em,
    itemsep=0.2em
}

\setlist[enumerate]{
    topsep=0.3em,
    itemsep=0.2em
}

\clubpenalty=10000
\widowpenalty=10000

% ============================================================
% ENCABEZADOS Y PIES
% ============================================================
\pagestyle{fancy}
\fancyhf{}

\fancyhead[L]{\nouppercase{\leftmark}}
\fancyhead[R]{\materia}

\fancyfoot[C]{\thepage}

\renewcommand{\headrulewidth}{0.4pt}
\renewcommand{\footrulewidth}{0pt}

\fancypagestyle{plain}{
    \fancyhf{}
    \fancyfoot[C]{\thepage}
    \renewcommand{\headrulewidth}{0pt}
}

% ============================================================
% COMANDOS MATEMÁTICOS
% ============================================================
\newcommand{\abs}[1]{\left|#1\right|}
\newcommand{\norm}[1]{\left\lVert#1\right\rVert}
\newcommand{\vect}[1]{\mathbf{#1}}
\newcommand{\deriv}[2]{\frac{d#1}{d#2}}
\newcommand{\pderiv}[2]{\frac{\partial#1}{\partial#2}}

% ============================================================
% CAJAS ACADÉMICAS
% ============================================================
\newtcolorbox{definition}[1]{
    enhanced,
    breakable,
    title={Definición: #1},
    fonttitle=\bfseries,
    colback=blue!3,
    colframe=blue!50!black,
    boxrule=0.8pt,
    arc=2mm
}

\newtcolorbox{important}{
    enhanced,
    breakable,
    title={Importante},
    fonttitle=\bfseries,
    colback=yellow!8,
    colframe=orange!70!black,
    boxrule=0.8pt,
    arc=2mm
}

\newtcolorbox{example}[1]{
    enhanced,
    breakable,
    title={Ejemplo: #1},
    fonttitle=\bfseries,
    colback=green!4,
    colframe=green!40!black,
    boxrule=0.8pt,
    arc=2mm
}

\newtcolorbox{formula}[1]{
    enhanced,
    breakable,
    title={Fórmula / Regla: #1},
    fonttitle=\bfseries,
    colback=purple!4,
    colframe=purple!50!black,
    boxrule=0.8pt,
    arc=2mm
}

\newtcolorbox{exercise}[1]{
    enhanced,
    breakable,
    title={Ejercicio: #1},
    fonttitle=\bfseries,
    colback=gray!5,
    colframe=gray!60!black,
    boxrule=0.8pt,
    arc=2mm
}

\newtcolorbox{note}{
    enhanced,
    breakable,
    title={Nota},
    fonttitle=\bfseries,
    colback=white,
    colframe=gray!50,
    boxrule=0.6pt,
    arc=2mm
}

% ============================================================
% HIPERVÍNCULOS Y METADATOS PDF
% ============================================================
\usepackage[
    colorlinks=true,
    linkcolor=blue!50!black,
    citecolor=green!40!black,
    urlcolor=blue!60!black,
    pdftitle={\titulo},
    pdfauthor={\autor},
    pdfsubject={Apuntes universitarios},
    bookmarks=true
]{hyperref}

% ============================================================
% DOCUMENTO
% ============================================================
\begin{document}

% ---------------- PORTADA ----------------
\begin{titlepage}

\thispagestyle{empty}

\begin{center}

\vspace*{1cm}

{\large \textsc{\universidad}}

\vspace{0.2cm}

{\large \carrera}

\vspace{3cm}

{\Huge \textbf{\materia}}

\vspace{1cm}

{\LARGE \titulo}

\vspace{0.8cm}

\textit{Comprender los conceptos en su idioma original es el primer paso para pensar el derecho más allá de las fronteras.}

\vspace{1.2cm}

\rule{0.70\textwidth}{0.5pt}

\vspace{0.5cm}

{\large Apuntes de estudio}

\vspace{0.5cm}

\rule{0.70\textwidth}{0.5pt}

\vfill

{\large \autor}

\vspace{0.3cm}

{\large \anio}

\end{center}

\vfill

\begin{flushleft}
\textit{Este es un modesto aporte a los estudiantes de ``Lectocomprensión'' no pretende sustituir el material oficial de la materia.}
\end{flushleft}

\end{titlepage}

% ---------------- ÍNDICE ----------------
\tableofcontents

% ============================================================
% CONTENIDO
% ============================================================
\chapter{Inglés jurídico: contratos y lectura de fallos}

La clase continúa el trabajo con textos de inglés jurídico. Se abordan dos textos: el primero, \textbf{``Offers and Invitations to Treat''}, distingue la oferta (\textbf{``offer''}), la contraoferta (\textbf{``counter offer''}) y la invitación a negociar (\textbf{``invitation to treat''}); el segundo, una síntesis del fallo \textbf{``Hammer v. Sidway''}, examina la contraprestación (\textbf{``consideration''}) y por qué abstenerse de una conducta puede constituirla. Asimismo, se compara la estructura de las sentencias en inglés con la de los fallos argentinos y se propone una estrategia de lectura basada en extraer ideas en lugar de traducir palabra por palabra.

\section*{Relevancia profesional de la lectocomprensión jurídica}
\addcontentsline{toc}{section}{Relevancia profesional de la lectocomprensión jurídica}

El dominio de la lectocomprensión en inglés jurídico resulta útil para el ejercicio profesional por razones concretas:

\begin{itemize}
    \item \textbf{Leer lo que otros no pueden leer.} La jurisprudencia comparada, los contratos internacionales y buena parte de la doctrina circulan en inglés. Saber dónde están los hechos, el derecho, el análisis y la decisión en una sentencia (\textbf{``facts''}, \textbf{``law''}, \textbf{``reasoning''}, \textbf{``order''}) ahorra horas de investigación.
    \item \textbf{Saber cuándo hay contrato y cuándo todavía no.} Distinguir \textbf{``invitation to treat''}, \textbf{``offer''} y \textbf{``counter offer''} permite detectar, en un mail, una propuesta o una publicación, si ya existe una obligación o si todavía se está negociando: evita comprometerse por accidente o reclamar algo que nunca existió.
    \item \textbf{Entender los contratos de todos los días.} Al aceptar términos online sin pagar, muchas veces se «paga» renunciando a derechos (publicidad, cookies, datos): eso es \textbf{``consideration''} en acción, aplicada a la vida cotidiana.
    \item \textbf{Leer mejor en cualquier idioma.} La estrategia de ir de la oración al párrafo, sin traducir cada palabra, sirve para exámenes, para el trabajo y para cualquier texto técnico.
\end{itemize}

\begin{example}{Cómo se relacionan los temas --- la compra de una bicicleta usada}
Al querer comprar una bicicleta usada en una plataforma de compraventa online, la publicación que se observa constituye una \textbf{``invitation to treat''}: invita a negociar, pero todavía no hay contrato. Cuando el comprador propone un precio, eso es una \textbf{``offer''}. Si el vendedor responde «te lo dejo a tanto», se produce una \textbf{``counter offer''}, y recién cuando una de las partes acepta sin cambios queda formado el contrato.

Para que ese contrato sea válido, cada parte debe dar algo a cambio: la \textbf{``consideration''}. En \textbf{``Hammer v. Sidway''} se discutió precisamente si una promesa a cambio de no hacer algo alcanzaba para constituirla. Si hay conflicto, un tribunal lo resuelve mediante un fallo, y conocer su estructura permite encontrar rápidamente los hechos, el derecho y la decisión. Para leer ese fallo en inglés jurídico no hace falta traducir cada palabra: alcanza con dominar el vocabulario clave y aplicar una buena estrategia de lectura.
\end{example}

\section{Anticipación del texto mediante la macroestructura}

Antes de iniciar la lectura del segundo texto de la unidad, \textbf{``Offers and Invitations to Treat''}, se recomienda observar su macroestructura: el gráfico en el que aparecen los términos \textbf{``offer''} y \textbf{``acceptance''}, ya trabajados en clases anteriores, y los subtítulos \textbf{``What is an offer?''} y \textbf{``What is an invitation to treat?''}. La presencia de estos subtítulos permite anticipar que cada apartado del texto ofrecerá una definición o conceptualización del término correspondiente, antes incluso de comenzar la lectura detallada.

\section{\texorpdfstring{``Offer'', ``counter offer'' e ``invitation to treat''}{Offer, counter offer e invitation to treat}}

\subsection{La ``offer'' como invitación a celebrar un contrato vinculante}

Una \textbf{``offer''} es, en primer lugar, una invitación de una persona a otra. Se retoma aquí el elemento de \textbf{``offer and acceptance''} ya estudiado como uno de los elementos constitutivos del contrato: debe existir esta aceptación para que el contrato exista.

La finalidad de la \textbf{``offer''} es celebrar un \textbf{``binding agreement''}, es decir, un contrato vinculante, sujeto a \textbf{``specific terms''} (condiciones particulares). La palabra \textbf{``binding''} no es azarosa: existen invitaciones dirigidas a otro tipo de pactos o acuerdos que no necesariamente tienen por finalidad celebrar un contrato con los efectos jurídicos propios de este. Cuando se produce la aceptación, el contrato queda conformado.

\begin{note}
\textbf{``...binding agreement... with specific terms''}

\textbf{``...the acceptance must result in a valid contract and is legally binding on the parties.''}
\end{note}

Debe existir el acto de aceptación para que se entienda que existe la voluntad de contratar.

\subsection{La definición legal de ``offer''}

El texto ofrece, además, una definición legal de \textbf{``offer''}, construida con sinónimos: \textbf{``proposal''}, \textbf{``a promise''} u \textbf{``other manifestation of willingness''}.

\begin{definition}{``Offer''}
\textbf{``...a proposal, a promise, or other manifestation of willingness to enter into a contract or bargain under proposed terms with another party that has the power to accept it.''}
\end{definition}

\textbf{``Willingness'' y ``will'': la idea de voluntad.} El sufijo \textbf{``-ness''} se agrega a numerosas palabras para sustantivarlas, es decir, para convertirlas en un sustantivo abstracto; de este modo, \textbf{``willingness''} significa voluntad. \textbf{``Will''} también significa voluntad y, en materia sucesoria, designa el testamento (\textbf{``last will and testament''}).

Esta idea de voluntad no es accesoria: existe una intención, una intencionalidad, de celebrar un contrato \textbf{``or bargain under proposed terms with another party that has the power to accept it''}. Es importante notar que esta noción no coincide con la idea de negociación --- el ida y vuelta previo a la aceptación ---, sino que describe el acto por el cual una parte declara su voluntad de contratar y la otra parte, que tiene el poder de aceptar, celebra el acto que termina de conformar dicha voluntad.

\subsection{\texorpdfstring{``Counter offer'': la aceptación condicionada}{Counter offer: la aceptación condicionada}}

Cuando la aceptación de una oferta se encuentra condicionada --- es decir, cuando introduce cambios en las condiciones propuestas ---, lo que se produce no es la celebración del contrato sino una \textbf{``counter offer''} (contraoferta).

\begin{example}{La contraoferta cotidiana}
Situación típica de una plataforma de compraventa online: ante la pregunta «¿me lo dejás a tanto?», no se está aceptando la oferta original, sino formulando una contraoferta. Quien efectivamente acepta --- si lo hace --- es la otra parte, en caso de aceptar ese cambio en las condiciones.
\end{example}

\subsection{\texorpdfstring{``Invitation to treat'': la invitación a negociar}{Invitation to treat: la invitación a negociar}}

La \textbf{``invitation to treat''} es una figura distinta de la \textbf{``offer''}, aunque relacionada con ella: constituye un paso previo a la oferta, es decir, la creación de las condiciones para que una oferta pueda existir.

\begin{definition}{``Invitation to treat''}
\textbf{``An invitation to treat is an invitation to start negotiations with intent to create an offer.''}

\textbf{``Definition: a solicitation for more offers, usually as a preliminary step to form a contract.''}
\end{definition}

La \textbf{``invitation to treat''} puede describirse como «abrir el juego»: una manera de decir «¿qué me proponés?», a partir de la cual recién podrá surgir la oferta propiamente dicha. En este momento todavía no existe la oferta como elemento formal del contrato: se trata de una propuesta genérica, no de un pedido con efectos jurídicos vinculantes.

\begin{example}{La reunión previa a la propuesta}
En contratos comerciales, una reunión virtual con el departamento de legales de otra empresa --- destinada a poner en contexto qué se quiere hacer y qué se necesita --- concluye típicamente en el pedido: «hacéme llegar una propuesta». Esa propuesta será la \textbf{``offer''}; la reunión previa es, en cambio, una instancia mucho más vaga desde el punto de vista jurídico, propia de la \textbf{``invitation to treat''}.
\end{example}

\subsection{Las tratativas preliminares en el derecho argentino}

La \textbf{``invitation to treat''} guarda cierta semejanza con las tratativas preliminares del derecho argentino. Se trata, sin embargo, de una comparación aproximada: el tema se retoma con mayor profundidad al abordar la jurisprudencia y el planteo de casos.

\subsection{\texorpdfstring{Diferencia entre ``offer'' e ``invitation to treat''}{Diferencia entre offer e invitation to treat}}

El texto plantea, en su último párrafo, la diferencia entre ambas figuras: el momento procesal de cada una es totalmente distinto. Lo que debe aceptarse es la oferta tal cual fue propuesta; si se introduce algún cambio, se está frente a una \textbf{``counter offer''}, que deja sin efecto la oferta original y da lugar a una nueva oferta que la otra parte deberá aceptar.

La \textbf{``invitation to treat''} se sitúa en un momento anterior: es la etapa de negociaciones previas, la etapa precontractual. En el derecho argentino no se la considera un elemento constitutivo del contrato; en todo caso, se la evalúa cuando el contrato es objeto de discusión judicial --- por ejemplo, si existió algún vicio de la voluntad.

\begin{table}[H]
\centering
\caption{Cuadro comparativo: figuras precontractuales y contractuales}
\begin{tabularx}{\textwidth}{
    >{\raggedright\arraybackslash}p{0.22\textwidth}
    >{\raggedright\arraybackslash}X
    >{\raggedright\arraybackslash}X
}
\toprule
\textbf{Figura} & \textbf{Qué es} & \textbf{Efecto jurídico} \\
\midrule
\textbf{``Invitation to treat''} &
Invitación a iniciar negociaciones con la intención de llegar a una oferta; «abrir el juego»; \textbf{``solicitation for more offers''}. &
Todavía no existe oferta ni contrato. Etapa precontractual; en el derecho argentino, algo parecido a las tratativas preliminares. Solo se evalúa si el contrato se discute en juicio. \\
\textbf{``Offer''} &
\textbf{``Proposal''}, \textbf{``promise''} u otra manifestación de voluntad (\textbf{``willingness''}) de celebrar un contrato bajo términos propuestos, dirigida a otra parte que tiene el poder de aceptarla. &
Elemento constitutivo del contrato (junto con la \textbf{``acceptance''}). Debe aceptarse tal cual fue propuesta. \\
\textbf{``Counter offer''} &
Aceptación condicionada o con cambios en las condiciones; típico «¿me lo dejás a tanto?». &
No es aceptación: la oferta original queda sin efecto y surge una nueva oferta que debe aceptar la otra parte. \\
\textbf{``Acceptance''} &
Acto por el cual se acepta la oferta \textbf{``on the parties''}. &
Debe existir para que exista la voluntad; produce un contrato válido y vinculante (\textbf{``legally binding''}). \\
\bottomrule
\end{tabularx}
\end{table}

\begin{important}
La \textbf{``invitation to treat''} está antes de la oferta: es la etapa de negociaciones previas. La \textbf{``offer''} aceptada tal cual forma el contrato; si la aceptación introduce cambios, hay una \textbf{``counter offer''} y comienza de nuevo el ciclo.
\end{important}

\section{Anticipación de un fallo: el texto 5}

\subsection{Presentación del texto}

El segundo texto trabajado en la clase, identificado como texto número cinco, se relaciona directamente con lo expuesto sobre \textbf{``offer''} e \textbf{``invitation to treat''}. Por su título, se anticipa que se trata de un fallo; más precisamente, de una síntesis de un fallo.

\subsection{Los títulos como adelanto del contenido}

Antes de la lectura detallada, los propios títulos del texto permiten anticipar su contenido. El primero, \textbf{``The issue of consideration''}, remite a la \textbf{``consideration''} --- la contraprestación, o las prestaciones recíprocas ---, uno de los elementos del contrato, lo que indica que este elemento estuvo en discusión en el fallo. El segundo, \textbf{``The conclusion''}, remite a la parte dispositiva o resolutiva, es decir, a lo que resolvió el tribunal. El tercero, \textbf{``Application to today's world''}, corresponde a un comentario sobre la relevancia actual del fallo: por qué importa analizarlo hoy.

\begin{table}[H]
\centering
\caption{Los títulos del texto 5 y lo que anticipan}
\begin{tabularx}{\textwidth}{
    >{\raggedright\arraybackslash}p{0.35\textwidth}
    >{\raggedright\arraybackslash}X
}
\toprule
\textbf{Título en inglés} & \textbf{Qué anticipa} \\
\midrule
\textbf{``The issue of consideration''} & La cuestión en debate: la contraprestación (\textbf{``consideration''}), uno de los elementos del contrato. \\
\textbf{``The conclusion''} & La parte dispositiva o resolutiva: lo que resolvió el tribunal. \\
\textbf{``Application to today's world''} & Un comentario sobre la relevancia actual del fallo: por qué importa analizarlo hoy. \\
\bottomrule
\end{tabularx}
\end{table}

\section{\texorpdfstring{Análisis del fallo ``Hammer v. Sidway''}{Análisis del fallo Hammer v. Sidway}}

\subsection{Los hechos del caso}

El fallo trata de un contrato entre un tío y su sobrino, que contiene una obligación de no hacer: una obligación negativa por la cual el sobrino se compromete a no incurrir en ciertos vicios (tomar, fumar, entre otras conductas asociadas a «la mala vida»). A cambio de que el sobrino se abstuviera de esas conductas, el tío le prometía una suma de dinero, exigible cuando el sobrino alcanzara la mayoría de edad, a los veintiún años.

% TODO: contenido incompleto o ambiguo en las notas originales
El monto prometido no queda del todo preciso en el registro de la clase: se menciona la cifra de «cinco mil» y, a continuación, de manera parcialmente inaudible, «setenta y cinco mil», en referencia aparente a una equivalencia en dinero actual.

\begin{example}{El tío, el sobrino y la promesa}
Un tío promete a su sobrino una suma de dinero si, hasta que cumpla la mayoría de edad (veintiún años), se abstiene de ciertas conductas: tomar, fumar y otros vicios (obligación de no hacer). Cuando el sobrino cumple los veintiún años, le pregunta al tío si recuerda la promesa que le hizo; el tío responde que sí, que le va a depositar el dinero. Pocos días después, el tío le escribe una carta al sobrino diciéndole que el dinero está en el banco y que puede considerarlo como dinero que genera intereses. El tío fallece; tenía además contratos con terceras personas que debían ejecutarse, y el pago al sobrino queda pendiente.
\end{example}

\subsection{Las partes: Hammer y Sidway}

El nombre del fallo, \textbf{``Hammer versus Sidway''}, no incluye a ninguna de las dos partes originales del contrato: el tío y el sobrino comparten el mismo nombre, William Story --- Senior y Junior (``the Second'') ---. Sidway es el \textbf{``executor''} del \textbf{``estate''} del tío, es decir, el albacea que ejecuta las disposiciones de última voluntad contenidas en el testamento (\textbf{``will''}); es la parte demandada.

Hammer, en cambio, es un tercero acreedor: quien reclama el pago es en verdad un cesionario, a quien el sobrino --- cansado de esperar --- cedió su crédito. Hammer es, por lo tanto, el actor, motivo por el cual su nombre encabeza la causa. El albacea se niega a pagar sosteniendo que el contrato entre el tío y el sobrino era inválido, porque el tío no había recibido nada a cambio: el sobrino solo había ofrecido la omisión de una acción, lo cual --- a criterio del albacea --- no constituía una contraprestación (\textbf{``consideration''}) válida.

\begin{table}[H]
\centering
\caption{\texorpdfstring{Partes y roles en ``Hammer v. Sidway''}{Partes y roles en Hammer v. Sidway}}
\begin{tabularx}{\textwidth}{
    >{\raggedright\arraybackslash}p{0.24\textwidth}
    >{\raggedright\arraybackslash}p{0.20\textwidth}
    >{\raggedright\arraybackslash}X
}
\toprule
\textbf{Sujeto} & \textbf{Rol} & \textbf{Explicación} \\
\midrule
Tío (William Story, Senior) & Promitente (\textbf{``promisor''}) & Promete el dinero a cambio de que el sobrino se abstenga de ciertos vicios. Fallece. \\
Sobrino (William Story, Junior, ``the Second'') & Beneficiario (\textbf{``promisee''}) & Cumple la abstención y, cansado de esperar, cede su crédito a un tercero. \\
Sidway & Albacea (\textbf{``executor''}) & Ejecuta las disposiciones de última voluntad; se niega a pagar. Es el demandado. \\
Hammer & Tercero acreedor, cesionario (\textbf{``a third party to whom the promise had been assigned''}) & Reclama el pago al albacea. Es el actor: por eso su nombre encabeza la causa. \\
\bottomrule
\end{tabularx}
\end{table}

\subsection{\texorpdfstring{La cuestión jurídica: la ``consideration'' y la obligación de no hacer}{La cuestión jurídica: la consideration y la obligación de no hacer}}

El punto discutido por el tribunal consistía en determinar si el hecho de no hacer --- de suspender una determinada conducta --- constituía una contraprestación suficiente para la existencia de un contrato válido.

\begin{important}
¿Puede constituir \textbf{``consideration''} (contraprestación válida) el hecho de abstenerse de una conducta que uno tenía derecho a realizar? El albacea sostenía que no: el tío daba dinero, pero el sobrino no daba nada; por lo tanto, faltaba un elemento y el contrato era inválido.
\end{important}

\subsection{La decisión: primera instancia y Cámara de Apelaciones}

En primera instancia se falló a favor del albacea: se consideró que no existía contraprestación válida y que, por ende, no había contrato. Sin embargo, el fallo fue apelado, y la Cámara de Apelaciones de Nueva York lo revocó, fallando a favor del tercero acreedor. El fundamento de la Cámara fue que el acto de abstenerse de realizar determinadas conductas constituye, en sí mismo, una contraprestación: esta no necesariamente debe consistir en una acción o en la entrega de algo, sino que también puede consistir en el cese o la suspensión de un derecho.

El fallo agrega que \textbf{``we need not speculate''}: no es necesario analizar cuánto le costó al sobrino dejar de fumar, de salir o de jugar. Basta con advertir que tenía un derecho legítimo a realizar esas conductas y que dejó de ejercerlo porque le habían prometido algo a cambio.

\begin{table}[H]
\centering
\caption{Recorrido procesal del caso}
\begin{tabularx}{\textwidth}{
    >{\raggedright\arraybackslash}p{0.26\textwidth}
    >{\raggedright\arraybackslash}p{0.28\textwidth}
    >{\raggedright\arraybackslash}X
}
\toprule
\textbf{Instancia} & \textbf{Resultado} & \textbf{Fundamento} \\
\midrule
Primera instancia & A favor del albacea (Sidway): no debe pagar. & No hay contraprestación válida; por ende, no hay contrato. \\
Cámara de Apelaciones de Nueva York & Revoca; falla a favor del tercero acreedor (Hammer). & La abstención de una conducta constituye contraprestación; no importa cuánto le costó al sobrino: renunció a un derecho legítimo porque se le prometió algo a cambio. \\
\bottomrule
\end{tabularx}
\end{table}

\subsection{\texorpdfstring{Vigencia del fallo: ``Application to today's world''}{Vigencia del fallo: Application to today's world}}

El texto destaca que los principios sentados por la Cámara resultan aplicables a los contratos que se forman en la actualidad, incluidos los contratos celebrados en línea.

\begin{example}{Los contratos online}
Al aceptar los términos de un servicio online sin pagar, se puede estar renunciando a un derecho --- navegar sin publicidad, sin que se guarden las cookies y los datos de navegación ---: esa renuncia funciona como la contraprestación (\textbf{``consideration''}) del acuerdo.
\end{example}

En definitiva, lo que constituye un contrato válido y vinculante (\textbf{``a valid contract''}) sigue girando en torno a la contraprestación: hasta dónde llega la prestación de cada parte, cuestión que en este caso estuvo directamente en tela de juicio.

\section{Estructura de las sentencias en inglés}

\subsection{Por qué conviene conocer la estructura}

Aunque el texto cinco no es el fallo en sí sino un comentario sobre él, resulta útil porque introduce la estructura habitual de los fallos en inglés, un aspecto con el que puede llegar a encontrarse el estudiante en el ejercicio de la profesión. Al igual que en los fallos argentinos --- con sus considerandos y votos ---, toda sentencia judicial comienza habitualmente con el \textbf{``case title''} y la mención de los jueces intervinientes.

\subsection{\texorpdfstring{``Case title'' y encabezado}{Case title y encabezado}}

En Estados Unidos, el encabezado de la sentencia suele presentarse dentro de un rectángulo, donde se indica el tribunal y las partes («fulano de tal, actora; fulano de tal, demandado»), con forma de cuadro. Esta presentación resulta visualmente más clara que la utilizada en los escritos judiciales argentinos, donde el encabezado suele ser todo texto corrido, lo que dificulta encontrar la información con rapidez.

\subsection{\texorpdfstring{``Findings of fact'': los hechos}{Findings of fact: los hechos}}

A continuación de las cuestiones de los hechos (\textbf{``facts''}), a veces referidas como \textbf{``findings of fact''}, se expone la determinación de los hechos que el tribunal considera pertinentes para la causa: toda la narrativa de los hechos. En el caso analizado, se relata que el tío prometió pagarle al sobrino, y que la cuestión de fondo --- lo que estaba en tela de juicio --- era si la omisión de actuar de determinada manera constituía una contraprestación suficiente. Para resolver esa cuestión de fondo se recurre a las \textbf{``determinations of law''}: el derecho aplicable.

\subsection{\texorpdfstring{``Determinations of law'': el derecho aplicable}{Determinations of law: el derecho aplicable}}

A diferencia de los fallos argentinos, donde los hechos y el derecho suelen aparecer mezclados --- lo que exige buscar la información interesante en el medio del texto ---, en las sentencias en inglés lo fáctico y lo legal aparecen claramente separados y marcados, lo que permite encontrar la información con mayor rapidez.

Las \textbf{``determinations of law''} refieren a cuál es la ley aplicable, entendida en sentido amplio: no solamente la ley escrita o positiva, emanada del Congreso, sino también los antecedentes jurisprudenciales, los principios jurídicos y, eventualmente, las buenas costumbres y los usos, es decir, el derecho en sentido amplio.

\subsection{\texorpdfstring{``Application'' y ``reasoning'': el análisis del tribunal}{Application y reasoning: el análisis del tribunal}}

Luego de establecer los hechos y el derecho aplicable, suele venir una parte dedicada a la aplicación de ese derecho, es decir, a cómo el tribunal conjuga ambos elementos. A esta parte se la denomina, indistintamente, \textbf{``application''} o \textbf{``reasoning''}. Este \textbf{``reasoning''} resulta comparable a los votos de los jueces en el sistema argentino: cada uno expone su interpretación del derecho aplicado a los hechos, lo que puede dar lugar a posturas diferentes entre los miembros de un tribunal colegiado antes de arribar a una resolución conjunta.

\subsection{\texorpdfstring{``Conclusions'', ``decision'', ``order'': la parte resolutiva}{Conclusions, decision, order: la parte resolutiva}}

Por último, del \textbf{``reasoning''} se desprenden las \textbf{``conclusions''}, la \textbf{``decision''} o la \textbf{``order''} --- distintas palabras que pueden aparecer para referirse a la parte dispositiva o resolutiva del fallo. Conocer estos términos resulta útil para la investigación jurídica: permite localizar con rapidez, dentro de un fallo extenso redactado en otra jurisdicción, la información puntual que se busca, sin necesidad de leer la totalidad del texto.

\begin{table}[H]
\centering
\caption{Estructura de una sentencia en inglés y en los fallos argentinos}
\begin{tabularx}{\textwidth}{
    >{\raggedright\arraybackslash}p{0.28\textwidth}
    >{\raggedright\arraybackslash}X
    >{\raggedright\arraybackslash}p{0.24\textwidth}
}
\toprule
\textbf{Parte en inglés} & \textbf{Contenido} & \textbf{En los fallos argentinos} \\
\midrule
\textbf{``Case title''} / encabezado & Partes, tribunal y jueces que intervienen; se presenta en un rectángulo. & Encabezado todo en texto; encontrar la información cuesta más. \\
\textbf{``Findings of fact''} / \textbf{``facts''} & Determinación y narrativa de los hechos que el tribunal considera pertinentes. & Los hechos y el derecho aparecen algo mezclados. \\
\textbf{``Determination(s) of law''} & Derecho aplicable en sentido amplio: normas, antecedentes jurisprudenciales, principios jurídicos, usos y costumbres. & Considerandos, donde lo jurídico y lo fáctico se entrelazan. \\
\textbf{``Application''} / \textbf{``reasoning''} & Aplicación del derecho a los hechos; cómo el tribunal conjuga ambos. & Parecido a los votos de cada juez. \\
\textbf{``Conclusions''} / \textbf{``decision''} / \textbf{``order''} & Parte dispositiva o resolutiva. & Parte dispositiva o resolución. \\
\bottomrule
\end{tabularx}
\end{table}

\begin{note}
Se trata del mismo tipo textual (sentencias judiciales), pero estructurado de otra manera. La información se encuentra con mayor facilidad en un fallo en inglés que en los fallos argentinos: el encabezado en rectángulo y la separación entre hechos, derecho, análisis y resolución permiten ubicar los datos con rapidez.
\end{note}

\section{Vocabulario y lectura detallada del texto 5}

\subsection{\texorpdfstring{La carta del tío, el ``estate'' y el ``executor''}{La carta del tío, el estate y el executor}}

El texto señala que, hacia el final del primer párrafo, \textbf{``a few days later, the uncle wrote the nephew a letter [...] and said that the money was in the bank and that he could consider this money on interest''}. Es decir: el tío no entregó el dinero de inmediato, sino que informó al sobrino que el dinero se encontraba depositado en el banco, generando intereses.

\begin{definition}{``Estate''}
El \textbf{``estate''} no es el testamento --- que se designa con la palabra \textbf{``will''} ---, sino el patrimonio o acervo sucesorio, es decir, aquello que el testamento distribuye o asigna. El testamento es el instrumento que reparte el \textbf{``estate''}.
\end{definition}

\subsection{\texorpdfstring{``Assigned'': la cesión del crédito}{Assigned: la cesión del crédito}}

El texto indica que \textbf{``the executor of the uncle's estate refused [...] interest claim brought by Hammer''}, aclarando luego que Hammer es \textbf{``a third party to whom the promise had been assigned''}: es decir, un cesionario del crédito que originalmente correspondía al sobrino, quien --- cansado de esperar --- decidió cederlo.

\begin{note}
No se cede «el contrato» en sí, sino los derechos (o las obligaciones) que surgen de él, como quien cede el derecho a cobrar un bono. En el caso, el sobrino cede a Hammer su derecho a cobrar la promesa del tío: Hammer es cesionario de ese crédito.
\end{note}

\subsection{\texorpdfstring{``Contend'' y el estado del derecho en 1891}{Contend y el estado del derecho en 1891}}

El verbo \textbf{``contend''} significa discutir, refutar o debatir. El texto indica que \textbf{``the uncle received nothing and [the nephew benefited by fulfilling the uncle's wishes]''}: el sobrino mejoró su situación al cumplir la voluntad del tío, mientras que este, a cambio, no había recibido nada material.

El fallo debe situarse temporalmente: el tío falleció en 1887 y el fallo es de 1891. En ese momento, la noción de \textbf{``consideration''} todavía estaba en discusión y se entendía, en general, como el intercambio de cosas materiales o la realización de una conducta positiva; estaba abierta la pregunta de si una obligación negativa --- una abstención, como en este caso --- podía constituir contraprestación. Por esta razón, el fallo constituye un \textbf{``leading case''}.

\begin{definition}{``Leading case''}
Un \textbf{``leading case''} es un caso que sienta jurisprudencia y abre el camino interpretativo. \textbf{``To lead''} significa guiar, marcar. A partir de \textbf{``Hammer v. Sidway''} se comenzó a interpretar que la \textbf{``consideration''} también puede consistir en un no hacer, es decir, en una omisión o abstención de conducta.
\end{definition}

\subsection{\texorpdfstring{``Plaintiff'', ``defendant'' y ``prosecutor''}{Plaintiff, defendant y prosecutor}}

La lectura del fallo exige distinguir con precisión los roles procesales, dado que el tío y el sobrino comparten el mismo nombre (William Story) y Hammer --- el tercero cesionario --- no se identifica hasta avanzada la narración de los hechos. Quien lleva adelante la acción es la parte actora; la otra parte es la demandada.

\begin{table}[H]
\centering
\caption{Roles procesales en inglés}
\begin{tabularx}{\textwidth}{
    >{\raggedright\arraybackslash}p{0.20\textwidth}
    >{\raggedright\arraybackslash}X
    >{\raggedright\arraybackslash}p{0.26\textwidth}
}
\toprule
\textbf{Término} & \textbf{Significado} & \textbf{En el caso analizado} \\
\midrule
\textbf{``Plaintiff''} & Actora, demandante; quien lleva a cabo la acción (el accionante). & En este caso civil, Hammer. \\
\textbf{``Defendant''} & Demandado; en materia penal, acusado o imputado según el momento del proceso. & En este caso, el albacea (Sidway). \\
\textbf{``Prosecutor''} & El fiscal: en materia penal no hay un particular que demande, sino el Estado. & No interviene en este caso, que es de naturaleza civil. \\
\bottomrule
\end{tabularx}
\end{table}

\subsection{\texorpdfstring{``Forbearance'' y ``upon the faith''}{Forbearance y upon the faith}}

\begin{definition}{``Forbearance''}
Abstención o autorrestricción: dejar de hacer algo que legítimamente se podría hacer, a cambio de un bien mayor. Es equivalente, en este contexto, a la idea de \textbf{``refrain''}.
\end{definition}

El texto agrega, en el mismo párrafo, que \textbf{``it is sufficient that he restricted his lawful freedom of action within certain prescribed limits upon the faith of his uncle's agreement''}. Este pasaje resulta especialmente relevante: indica que el sobrino tenía un derecho legítimo, una libertad legítima, y que restringió esa libertad \textbf{``upon the faith''} --- es decir, confiando en la promesa de su tío. De no haber existido esa promesa, presumiblemente el sobrino habría continuado con las conductas de las que se abstuvo.

\begin{note}
\textbf{``Forbearance''}: abstención; autorrestricción; renunciar a algo que legítimamente se podía hacer, a cambio de un bien mayor. \textbf{``Upon the faith of''}: confiando en la promesa. La \textbf{``consideration''} del sobrino fue restringir su libertad legítima confiando en el acuerdo con su tío.
\end{note}

\subsection{\texorpdfstring{``Indulge'': darse el gusto}{Indulge: darse el gusto}}

El verbo \textbf{``indulge''}, presente en el primer párrafo del fallo, no equivale a «recaer», sino a permitirse o darse el gusto de algo: en el caso, se refiere a que el sobrino se permitía ciertos vicios. Se trata de una forma relativamente cuidada de indicar que el sobrino incurría en conductas lícitas --- tomar, fumar, jugar --- pero no ilícitas, en la medida en que no cometiera ningún delito. Precisamente porque esas conductas eran lícitas, la discusión jurídica giró en torno a si abstenerse de ellas podía constituir una contraprestación válida.

\subsection{\texorpdfstring{``Promisor'' y ``promisee''}{Promisor y promisee}}

Las partes de la promesa se identifican, en inglés jurídico, mediante las terminaciones \textbf{``-or''} y \textbf{``-ee''}: la primera para quien inicia o dirige el acto jurídico, y la segunda para quien lo recibe. Esta estructura se repite en buena parte de los negocios jurídicos en inglés.

\begin{table}[H]
\centering
\caption{Terminaciones ``-or'' y ``-ee'' en los negocios jurídicos}
\begin{tabularx}{\textwidth}{
    >{\raggedright\arraybackslash}p{0.20\textwidth}
    >{\raggedright\arraybackslash}X
    >{\raggedright\arraybackslash}p{0.18\textwidth}
}
\toprule
\textbf{Término} & \textbf{Quién es} & \textbf{En el caso} \\
\midrule
\textbf{``Promisor''} & Quien emite o dirige la promesa; quien inicia el acto (terminación \textbf{``-or''}). & El tío. \\
\textbf{``Promisee''} & Beneficiario o receptor de la promesa (terminación \textbf{``-ee''}). & El sobrino. \\
\bottomrule
\end{tabularx}
\end{table}

No siempre existe en español una traducción literal para estos términos: \textbf{``promisee''}, por ejemplo, no se traduce naturalmente como «prometido».

\subsection{\texorpdfstring{``The issue before the court arose...''}{The issue before the court arose...}}

El verbo \textbf{``arose''} es el pasado de \textbf{``arise''} (surgir, emerger). La expresión \textbf{``the issue before the court arose''} se refiere, entonces, a la cuestión o el conflicto que surgió ante el tribunal. En el pasaje correspondiente se remarca que el sobrino \textbf{``had performed''}: había cumplido con su parte, por lo cual tenía derecho a reclamar lo prometido.

\section{Estrategia de lectura: de las unidades pequeñas a las grandes}

La estrategia de lectura propuesta para el inglés jurídico consiste en avanzar de las unidades pequeñas hacia las más grandes: primero extraer la idea de la oración, luego la del párrafo, sin detenerse en cada palabra individual. La mayoría de las veces es posible evitar recurrir al diccionario; corresponde hacerlo, en cambio, ante los términos técnicos que condicionan la comprensión global del texto --- por ejemplo, \textbf{``consideration''}, o la comprensión del rol de un tercero como Hammer ---, dado que un error en estos términos puede llevar a una comprensión errónea de todo el texto.

\begin{note}
Ir de las unidades pequeñas a las grandes: primero la idea de la oración, luego la del párrafo, sin detenerse en cada palabra. Abrir el diccionario solo para los términos técnicos que condicionan la comprensión (por ejemplo, \textbf{``consideration''}, o el rol de Hammer como tercero acreedor). Dejar de lado la necesidad de entenderlo todo y de traducir literalmente.
\end{note}

Esta propuesta metodológica es aplicable no solo al inglés jurídico, sino a la lectura de cualquier texto técnico: se trata de correr la vista del propósito de «entender todo, traducir todo literalmente» hacia la extracción de ideas generales.

% ============================================================
% CUADRO CORNELL
% ============================================================
\section{Cuadro Cornell de repaso}

\begin{longtable}{
    >{\raggedright\arraybackslash}p{0.30\textwidth}
    >{\raggedright\arraybackslash}p{0.60\textwidth}
}
\caption{Cuadro Cornell --- Inglés jurídico: contratos y lectura de fallos} \\
\toprule
\textbf{Preguntas / palabras clave} & \textbf{Notas / respuestas} \\
\midrule
\endfirsthead

\multicolumn{2}{c}{\tablename\ \thetable{} (continuación)} \\
\toprule
\textbf{Preguntas / palabras clave} & \textbf{Notas / respuestas} \\
\midrule
\endhead

\midrule
\multicolumn{2}{r}{Continúa en la página siguiente} \\
\endfoot

\bottomrule
\endlastfoot

\textbf{¿Qué es una ``offer''?} &
Invitación de una persona a otra para celebrar un contrato vinculante (\textbf{``binding agreement''}) bajo términos específicos. El texto la define con sinónimos: \textbf{``proposal''}, \textbf{``promise''} u otra manifestación de voluntad (\textbf{``willingness''}). Cuando la otra parte acepta, queda conformado el contrato: la aceptación debe dar como resultado un contrato válido y vinculante (\textbf{``legally binding on the parties''}). \\

\textbf{¿Qué significan ``willingness'' y ``will''?} &
El sufijo \textbf{``-ness''} sustantiva y vuelve abstracta la palabra: \textbf{``willingness''} es voluntad. \textbf{``Will''} también es voluntad y, en sucesiones, el testamento (\textbf{``last will and testament''}). La oferta expresa una intención de contratar. \\

\textbf{¿Qué es una ``counter offer''?} &
Una aceptación condicionada o con cambios: no es aceptación. La oferta original queda sin efecto y surge una nueva oferta (la contraoferta), que debe aceptar la otra parte. Ejemplo: «¿me lo dejás a tanto?» en una plataforma de compraventa online. \\

\textbf{¿Qué es una ``invitation to treat''?} &
Invitación a iniciar negociaciones con la intención de llegar a crear una oferta (\textbf{``solicitation for more offers''}, paso preliminar). Es «abrir el juego»: «¿qué me proponés?». Todavía no existe la oferta como elemento formal del contrato. \\

\textbf{¿Cómo se relaciona con el derecho argentino?} &
Se parece a las tratativas preliminares: etapa precontractual. No es constitutiva del contrato; solo se evalúa si el contrato se discute en juicio (por ejemplo, por vicios de la voluntad). \\

\textbf{¿Qué es la ``consideration''?} &
La contraprestación: uno de los elementos del contrato (prestaciones recíprocas). Sin ella, el contrato se considera inválido. \\

\textbf{¿Cuáles son los hechos de ``Hammer v. Sidway''?} &
Un tío promete a su sobrino una suma de dinero si, hasta los veintiún años, se abstiene de tomar, fumar y otros vicios. El sobrino cumple; el tío le informa que el dinero está en el banco, con intereses, y fallece. El sobrino cede su crédito a un tercero. \\

\textbf{¿Quiénes son Hammer y Sidway?} &
Hammer: tercero acreedor, cesionario del crédito del sobrino; es el actor. Sidway: el albacea (\textbf{``executor''}) del patrimonio (\textbf{``estate''}) del tío; es el demandado y se niega a pagar. \\

\textbf{¿Cuál fue el conflicto jurídico y cómo se resolvió?} &
Si abstenerse de una conducta lícita es contraprestación suficiente. En primera instancia se falló a favor del albacea (sin contraprestación no hay contrato). La Cámara de Apelaciones de Nueva York revocó: la abstención (\textbf{``forbearance''}) es contraprestación; no importa cuánto le costó al sobrino (\textbf{``we need not speculate''}). \\

\textbf{¿Por qué es un ``leading case'' y qué vigencia tiene?} &
Sentó jurisprudencia: desde 1891 se interpretó que la \textbf{``consideration''} también puede ser un no hacer. Sus principios se aplican a contratos actuales, incluidos los online, donde se paga con renuncias de derechos (publicidad, cookies, datos). \\

\textbf{¿Cómo se estructura una sentencia en inglés?} &
\textbf{``Case title''} (rectángulo con partes y tribunal) $\rightarrow$ \textbf{``findings of fact''} $\rightarrow$ \textbf{``determinations of law''} $\rightarrow$ \textbf{``application''} / \textbf{``analysis''} / \textbf{``reasoning''} $\rightarrow$ \textbf{``conclusions''} / \textbf{``decision''} / \textbf{``order''}. Los hechos y el derecho están separados y marcados. \\

\textbf{¿Qué diferencia hay con los fallos argentinos?} &
Es el mismo tipo textual, pero en los fallos argentinos el encabezado es texto corrido y los hechos y el derecho aparecen mezclados; la información cuesta más de encontrar. Las sentencias en inglés resultan, en general, más claras y fáciles de interpretar. \\

\textbf{¿Qué son ``plaintiff'', ``defendant'' y ``prosecutor''?} &
\textbf{``Plaintiff''}: actora/demandante. \textbf{``Defendant''}: demandado (y, en penal, acusado o imputado). \textbf{``Prosecutor''}: fiscal, cuando quien acciona es el Estado y no un particular. \\

\textbf{¿Qué son ``promisor'' y ``promisee''?} &
Terminación \textbf{``-or''} para quien inicia el acto (el tío, que promete) y \textbf{``-ee''} para quien lo recibe (el sobrino, beneficiario de la promesa). No siempre tienen traducción literal. \\

\textbf{Vocabulario clave del fallo} &
\textbf{``Estate''}: patrimonio o acervo sucesorio. \textbf{``Assigned''}: cedido (se ceden derechos, no «el contrato»). \textbf{``Contend''}: discutir, refutar. \textbf{``Forbearance''}: abstención, autorrestricción. \textbf{``Upon the faith''}: confiando en la promesa. \textbf{``Indulge''}: permitirse, darse el gusto. \textbf{``Arose''}: surgió. \\

\textbf{¿Cómo leer mejor en inglés jurídico?} &
De la oración al párrafo, sin detenerse en cada palabra. Usar el diccionario solo para términos técnicos que condicionan la comprensión. No buscar traducir todo literalmente. \\

\midrule
\multicolumn{2}{p{0.90\textwidth}}{
\textbf{Resumen:} La clase recorrió dos textos que se complementan. En \textbf{``Offers and Invitations to Treat''} se distinguieron tres momentos de la formación del contrato: la invitación a negociar (\textbf{``invitation to treat''}), previa y no constitutiva; la oferta (\textbf{``offer''}), manifestación de voluntad de contratar que debe aceptarse tal cual; y la contraoferta (\textbf{``counter offer''}), que aparece cuando la aceptación cambia las condiciones y reinicia el ciclo. Con \textbf{``Hammer v. Sidway''} se analizó un caso en el que la falta de \textbf{``consideration''} fue el argumento para negar el pago, y se estableció, como \textbf{``leading case''}, que la abstención de una conducta lícita es contraprestación válida; el caso exige además comprender la cesión de créditos, el rol del albacea y los roles de \textbf{``promisor''} y \textbf{``promisee''}. A partir del fallo se comparó la estructura de las sentencias en inglés (\textbf{``case title''}, \textbf{``findings of fact''}, \textbf{``determinations of law''}, \textbf{``reasoning''}, \textbf{``conclusions''}) con la de las sentencias argentinas, concluyendo que la primera permite ubicar la información con mayor rapidez. Por último, se propuso una estrategia de lectura: ir de las ideas pequeñas a las grandes, apoyarse en la macroestructura y detenerse solo en el vocabulario técnico imprescindible.
} \\

\end{longtable}

\end{document}
